\section{Introduction}
\label{sec:intro}

NDH is a geometric cognitive emitter defined over a modal-partitioned manifold
with a triadic recursion kernel, spectral basis, and non-dual algebra. Its
internal geometry—including modal boundaries, holonomy clamps, spectral
transport, and recursion preservation—has been formalized in prior work. What
remains uncharacterized is the geometry of \emph{interaction}: how semantic
inputs induce traversal of NDH's cognitive manifold, how spectral components
deform under excitation, and how holonomy return operators classify the
resulting behavior.

Traditional machine learning systems, particularly transformer-based
architectures, lack the geometric structures required to model interaction at
this level. Transformers operate in a flat vector space without curvature,
parallel transport, modal adjacency, or fiber bundle structure. As a result,
their interaction behavior is governed by algebraic softmax dynamics rather
than geometric transport. NDH, by contrast, contains a connection, modal
boundaries, spectral bundles, and holonomy constraints, making interaction
inherently geometric.

This preprint develops the interaction geometry of NDH. We introduce the
cognitive manifold $\mathcal{M}_{cog}$, a geometric space governing
interrogation of NDH. Semantic inputs are modeled as excitation vectors in
$T_x \mathcal{M}_{cog}$, generating geodesic trajectories whose modal-boundary
crossings determine structural response. We define spectral bundles whose
fibers encode spectral emission under interrogation, and show that semantic
excitation induces section deformation and modal twisting of these bundles.

Holonomy return operators associated with closed interaction loops yield a
geometric taxonomy of NDH interactions. These operators classify responses into
modal-aligned, spectral-aligned, recursion-aligned, and holonomy-deviant types.
This classification is geometric, not statistical, and arises from curvature,
modal adjacency, and spectral bundle morphisms.

The goal of this work is to provide a formal, geometric account of NDH
interaction behavior, complementing the runtime geometry established
previously. This establishes NDH as a coherent cognitive manifold with
well-defined interaction dynamics, enabling rigorous analysis of semantic
excitation, spectral transport, and holonomy-based classification.
